\section{Results}

This section presents and analyses the results obtained from the three implemented incremental training configurations. These configurations were designed to evaluate the model’s ability to retain prior knowledge, adapt to new tasks, and mitigate the effects of \textit{catastrophic forgetting}, a common phenomenon in continual learning scenarios.

\begin{itemize}
  \item \textbf{Baseline:} reference model without memory or retention mechanisms.
  \item \textbf{Dual-Bullinaria:} implementation of the approach proposed by Bullinaria (2009), with partial weight duplication.
  \item \textbf{Dual-Bullinaria ×2:} extended version of the Dual model with reinforced duplication to increase learning stability.
\end{itemize}

Each of these variants presents a particular architecture and a distinct approach to preserving acquired knowledge. Table~\ref{tab:perfiles_incrementales} summarises their main characteristics and experimental purposes.

\begin{table}[H]
  \centering
  \renewcommand{\arraystretch}{1.3}
  \setlength{\tabcolsep}{6pt}
  \begin{tabular}{|p{3cm}|p{7cm}|p{5cm}|}
    \hline
    \textbf{Profile} & \textbf{Description} & \textbf{Purpose} \\ \hline
    \textbf{Baseline} & Simple incremental training without explicit memory or retention mechanisms. & Serve as a reference point for the model’s standard behaviour. \\ \hline
    \textbf{Dual-Bullinaria} & Recreation of Bullinaria’s (2009) protocol, introducing partial weight duplication to preserve information learned in previous stages. & Evaluate the effectiveness of the duplication mechanism in reducing catastrophic forgetting. \\ \hline
    \textbf{Dual-Bullinaria ×2} & Extension of the Dual profile with reinforced weight duplication per training block, emulating a higher-capacity memory. & Increase knowledge retention and stability throughout incremental learning. \\ \hline
  \end{tabular}
  \caption{Description of the incremental training profiles.}
  \label{tab:perfiles_incrementales}
\end{table}

\subsection{Performance Analysis}

Figure~\ref{fig:catastrophic_forgetting} illustrates the performance curves obtained during the training and evaluation phases. These curves reflect the behaviour of each model as new information blocks are progressively introduced, showing to what extent they manage to preserve previously learned knowledge.

\begin{figure}[H]
  \centering
  \includegraphics[width=0.85\textwidth]{results/catastrophic_forgetting}
  \caption{Evidence of catastrophic forgetting across the evaluated profiles.}
  \label{fig:catastrophic_forgetting}
\end{figure}

According to the experimental results, the following general trends were identified:

\begin{enumerate}
  \item \textbf{Baseline:} shows a gradual decrease in accuracy as new tasks are incorporated, reflecting the absence of memory mechanisms and, consequently, a progressive loss of prior knowledge.
  \item \textbf{Dual-Bullinaria:} maintains a more stable performance curve, showing a considerable reduction in forgetting and better preservation of learned information.
  \item \textbf{Dual-Bullinaria ×2:} reinforces knowledge retention, although slight initial fluctuations are observed due to the adjustment of duplicated weights during the early stages of training.
\end{enumerate}

\begin{table}[H]
  \centering
  \renewcommand{\arraystretch}{1.3}
  \setlength{\tabcolsep}{6pt}
  \begin{tabular}{|p{3cm}|p{6cm}|p{6cm}|}
    \hline
    \textbf{Profile} & \textbf{Retention} & \textbf{Main Observation} \\ \hline
    \textbf{Baseline} & Slight drop in accuracy (moderate forgetting). & Lacks memory mechanisms, leading to cumulative information loss. \\ \hline
    \textbf{Dual-Bullinaria} & Stability across tasks. & Achieves a balance between incorporating new information and preserving previous knowledge. \\ \hline
    \textbf{Dual-Bullinaria ×2} & Initial fluctuations followed by greater stability. & Strengthens long-term memory, although slightly reducing model adaptability. \\ \hline
  \end{tabular}
  \caption{Summary of experimental observations.}
  \label{tab:observaciones_perfiles}
\end{table}

\subsection{General Profile Comparison}

Figure~\ref{fig:perfil_comparation} presents a comparative view of the three approaches, highlighting differences in accuracy, retention, and stability among the profiles. A progressive improvement in memory capacity is observed as weight duplication strategies are incorporated, albeit at the cost of increased computational expense.

\begin{figure}[H]
  \centering
  \includegraphics[width=0.85\textwidth]{results/perfil_comparation}
  \caption{General comparison between the evaluated profiles.}
  \label{fig:perfil_comparation}
\end{figure}

In general terms, the observations can be summarised as follows:

\begin{itemize}
  \item \textbf{Baseline:} simplest model, prone to forgetting and with low stability in sequential tasks.
  \item \textbf{Dual-Bullinaria:} achieves a proper balance between memory and accuracy, showing a significant improvement in retention without compromising efficiency.
  \item \textbf{Dual-Bullinaria ×2:} greatly increases memory capacity, though its adaptability decreases, likely due to overfitting in the early training phases.
\end{itemize}

\begin{table}[H]
  \centering
  \renewcommand{\arraystretch}{1.25}
  \begin{tabularx}{\textwidth}{p{2.6cm} X X p{3.2cm}}
    \toprule
    \textbf{Profile} & \textbf{Description} & \textbf{Observed Behaviour} & \textbf{Conclusion} \\
    \midrule
    \textbf{Baseline} & Incremental training without retention mechanisms. & Shows good initial performance but suffers significant knowledge loss as new tasks are added. & Serves as a baseline for comparing the effectiveness of the other strategies. \\
    \textbf{Dual-Bullinaria} & Partial weight duplication to preserve prior knowledge. & Maintains a stable curve in ``Prev Accuracy'' and sustained performance in ``Test Accuracy''. & Represents the best balance between memory, stability, and continual learning capability. \\
    \textbf{Dual-Bullinaria $\times$ 2} & Extended duplication reinforcing weight preservation. & Exhibits greater knowledge retention with slight initial variations. & Enhances long-term stability, although it requires greater adjustment and computational resources. \\
    \bottomrule
  \end{tabularx}
  \caption{General comparison of incremental learning profiles.}
  \label{tab:perfiles_comparation}
\end{table}

\subsection{Conclusions}

The experimental results demonstrate that introducing weight duplication mechanisms is an effective strategy to mitigate \textit{catastrophic forgetting}, promoting greater stability and knowledge retention throughout incremental learning. However, this improvement entails an increase in model complexity, both in hyperparameter tuning and in computational resource consumption.

In summary, the \textbf{Dual-Bullinaria} and \textbf{Dual-Bullinaria ×2} profiles validate the effectiveness of structured-memory-based strategies. The former provides an optimal balance between retention and adaptability, whereas the latter achieves superior retention at the expense of reduced flexibility in incorporating new patterns.
